\documentclass{beamer}
% \documentclass[handout,xcolor=pdftex,dvipsnames,table]{beamer}

\mode<presentation>
{
  \usetheme{Warsaw}
  \usecolortheme{whale}
  % or ...

  % or whatever (possibly just delete it)
  \setbeamertemplate{navigation symbols}{}
}


% make a command to decrease font size locally when desired
\newcommand\Fontvi{\fontsize{8.5}{7.2}\selectfont}



\setbeamertemplate{itemize items}[ball]
\setbeamertemplate{itemize subitem}[triangle]
\setbeamertemplate{itemize subsubitem}[circle]
\setbeamercovered{invisible}

\usepackage{palatino} 
\usepackage{listings} % Gives syntax highlighting for python code. 
\usepackage{color} % Used for syntax highlighting. 
\usepackage{textcomp} % Used for syntax highlighting. 
\usepackage{caption}
\captionsetup{labelformat=empty,labelsep=none}
% This gives syntax highlighting in the python environment 
\definecolor{gray}{gray}{0.5} 
\definecolor{key}{rgb}{0,0.5,0} 
\lstset{
language=python,
basicstyle=\ttfamily\tiny, 
otherkeywords={1, 2, 3, 4, 5, 6, 7, 8 ,9 , 0, -, =, +, [, ], (, ), \{, \}, :, *, !}, 
keywordstyle=\color{blue}, 
stringstyle=\color{red},
showstringspaces=false,
alsoletter={1234567890},
otherkeywords={\ , \}, \{},
keywordstyle=\color{blue},
emph={access,and,break,class,continue,def,del,elif ,else,%
except,exec,finally,for,from,global,if,import,in,is,%
lambda,not,or,pass,print,raise,return,try,while},
emphstyle=\color{black}\bfseries,
emph={[2]True, False, None, self},
emphstyle=[2]\color{green},
emph={[3]from, import, as},
emphstyle=[3]\color{blue},
upquote=true,
morecomment=[s]{"""}{"""},
commentstyle=\color{gray}\slshape,
emph={[4]1, 2, 3, 4, 5, 6, 7, 8, 9, 0},
emphstyle=[4]\color{blue},
literate=*{:}{{\textcolor{blue}:}}{1}%
{=}{{\textcolor{blue}=}}{1}%
{-}{{\textcolor{blue}-}}{1}%
{+}{{\textcolor{blue}+}}{1}%
{*}{{\textcolor{blue}*}}{1}%
{!}{{\textcolor{blue}!}}{1}%
{(}{{\textcolor{blue}(}}{1}%
{)}{{\textcolor{blue})}}{1}%
{[}{{\textcolor{blue}[}}{1}%
{]}{{\textcolor{blue}]}}{1}%
{<}{{\textcolor{blue}<}}{1}%
{>}{{\textcolor{blue}>}}{1},%
numbers=none,
}

\newcommand{\putat}[3]{\begin{picture}(0,0)(0,0)\put(#1,#2){#3}\end{picture}}



\title[]{Python Workshop\\
File Input/Output}

\author[Fienen] % (optional, use only with lots of authors)
{Mike~Fienen}
\institute[USGS] % (optional, but mostly needed)
{
  U.S. Geological Survey\\
  Wisconsin Water Science Center, Middleton, Wisconsin USA
  }
  \titlegraphic{\includegraphics[scale=0.5]{figures/c_USGSid1.pdf}}
  

\date[UQ12] % (optional, should be abbreviation of conference name)
{USGS National Groundwater Workshop, August 2012}

\subject{Python}


\begin{document}

\begin{frame}
  \titlepage
\end{frame}
\logo{\vspace{-0.3cm} \includegraphics[width=1.5cm]{figures/c_USGSid1.pdf}\hspace*{11.10cm}}  

\begin{frame}{Outline}
\tableofcontents
\end{frame}

\begin{frame}[fragile]
\frametitle{Overview}
\begin{itemize}

\item Much of what is useful to do in Python is reading files, manipulating the data, and writing out results in another format
\item Python and Numpy provide ways to read and write ASCII and binary files. We will focus on ASCII files
\end{itemize}
\end{frame}

\begin{frame}[fragile]
\frametitle{Background Information}
\url{http://waterservices.usgs.gov/}
\url{http://nwis.waterdata.usgs.gov/nwis/pmcodes/}
\url{http://docs.python.org/tutorial/inputoutput.html}
\end{frame}


\begin{frame}[fragile]
\frametitle{Reading and Writing with Strings}
\begin{itemize}
\item The simplest way to write information to a string is using \texttt{str}
\begin{lstlisting}
>>>a = 5.4
>>>str(a)
'5.4'
\end{lstlisting}
\item We typically want more control. Two main ways to do it. Old school (\texttt{\%}) and new school (\texttt{format})
\item Formatted input and output are a key difference between Python 2.X and 3.X

\end{itemize}
\end{frame}


\begin{frame}[fragile]
\frametitle{Writing Strings the Old School Way (\texttt{\%})}
\begin{itemize}
\item The general syntax is to make a string with conversion types for variables. For example:
\begin{lstlisting}
>>>outstr = 'I have %21.1f kg of %s and %d bins of %s' %(3.99983,'eggs',53,'spam')
>>> outstr
'I have                   4.0 kg of eggs and 53 bins of spam'
>>>outstr = 'I have %-21.1f kg of %s and %0.3d bins of %s' %(3.99983,'eggs',53,'spam')
>>> outstr
'I have 4.0                   kg of eggs and 053 bins of spam'

\end{lstlisting}
\item The general idea is to make a string including \texttt{'\%'}, a conversion flag (optional), a width and resolution (optional), and a conversion type (required).
\item[] For example: \texttt{\%<flag><width>.<resolution><type>}\newline{}\texttt{\%-12.3f} Is a left-justified, floating point value with width of 12 and 3 decimal places.
\item Following the format string must be a list of values as a tuple identified by \texttt{\%}
\item[] \texttt{'\%4d is the \%s\textbackslash{}n' \%(42,'answer')}
\end{itemize}
\end{frame}



\begin{frame}[fragile]
\frametitle{Writing Strings the Old School Way (\texttt{\%})}
Details about formatted output available at:\newline{}
\begin{small} \url{http://docs.python.org/library/stdtypes.html} \end{small}

\begin{itemize}
\item Conversion flag characters
\begin{description}
\Fontvi
\item [\texttt{'\#'}] Invokes alternate behavior (see website for details)
\item [\texttt{'0'}] Pads numeric values with zeros
\item [\texttt{'-'}] Left-adjusts the output
\item [\texttt{' '}] Leave a space before signed positive values so they line up with negative ones

\end{description}
\pause
\item Most common conversion types. 
\end{itemize}
\begin{description}
\Fontvi
\item [\texttt{\%d} or \texttt{\%i} ] Signed integer
\item [\texttt{\%f} or \texttt{\%F} ] Floating point
\item [\texttt{\%e} or \texttt{\%E} ] Floating point exponential (lower or upper case)
\item [\texttt{\%g} or \texttt{\%G} ] Combination of \texttt{\%f} and \texttt{\%e} depending on resolution
\item [\texttt{\%s} or \texttt{\%r} ] String. Width is used, but not resolution

\end{description}
\end{frame}

\begin{frame}[fragile]
\frametitle{Writing Strings the New School Way (\texttt{format})}
Details about new school string formatting at:\newline{}
\begin{small}
\Fontvi
\url{http://docs.python.org/library/string.html#formatstrings}
\end{small}
\begin{itemize}
\item The general syntax is similar, but conversion information is supplied differently. For example:
\begin{lstlisting}
>>>outstr = 'I have %{0:21.1f} kg of {1:s} and {2:0=3d} bins of {3:s}'.format(3.99983,'eggs',53,'spam')
>>> outstr
'I have                   4.0 kg of eggs and 053 bins of spam'
\end{lstlisting}
\item In this case, make a string including \texttt{\{...\}}, statements with conversion information.

The general pattern is \texttt{\{[index]:[format]\}}\newline{}
\begin{itemize}
\item The \texttt{[index]} argument refers to the item index being mapped in\newline{}
\item The \texttt{[format]} argument is similar to those in the old school way, but with some additional flexibility
\end{itemize}

\end{itemize}
\end{frame}

\begin{frame}[fragile]
Now that we can write strings, how about we write them to files? \newline{}
We should also know how to read stuff back in from files
\end{frame}

\begin{frame}[fragile]
\frametitle{Interacting with Text Files}
\begin{itemize}
\item The first thing is to open a file and make a file object
\begin{lstlisting}
ifp = open('somefile.txt','r')
ofp = open('someotherfile.txt','w')
\end{lstlisting}
\begin{itemize}
\item This object can be used to read or write from. \\I use \texttt{ifp} for ``input file pointer" and \texttt{ofp} for ``output file pointer" \\
The arguments \texttt{'r'} and \texttt{'w'} indicate ``read" and ``write" respectively.
\end{itemize}
\item To read the file can use \texttt{readline()} or \texttt{readlines()}\\
\begin{itemize}
\item The difference is that  \texttt{readlines()} reads the entire file into memory rather than  \texttt{readline()}  which reads one line at a time. Most of the time,  \texttt{readlines()}  is better
\item With \texttt{readlines()}  once the data are read in, the result is a list with each element representing a line in the text file
  
\end{itemize}
\end{itemize}
\end{frame}


\begin{frame}[fragile]
\frametitle{\texttt{np.genfromtxt}: flexible way to read columns}
Example file: STATE\_FIPS.csv \\
\texttt{State Abbreviation,FIPS Code,State Name} \\
\texttt{AK,02,ALASKA} \\
\texttt{AL,01,ALABAMA}  \\
\pause
\begin{lstlisting}
import numpy as np
infilename = 'STATE_FIPS.csv'
indat = np.genfromtxt(infilename,delimiter=',',dtype=None,names=True)
\end{lstlisting}
\pause
\begin{description}
\item[\texttt{delimiter=','}] delimiter can be \emph{anything}
\item[\texttt{dtype=None}] Numpy interprets column data types. If unknown, makes it a string
\item[\texttt{names=True}] Each column gets a data type and a name
\end{description}
\pause
\begin{lstlisting}
In [10]: indat
Out[7]: 
array([('AK', 2, 'ALASKA'), ('AL', 1, 'ALABAMA'), ('AR', 5, 'ARKANSAS'), ... 
dtype=[('State_Abbreviation', '|S2'), ('FIPS_Code', '<i4'), ('State_Name', '|S20')])
\end{lstlisting}
\pause
\hskip 1.5 cm \emph{N.B.$\rightarrow$underscore replaces space in names}
\end{frame}


\end{document}